% Exported from Math++. Compile with pdflatex.
\documentclass[11pt,a4paper]{article}
\usepackage[T1]{fontenc}
\usepackage[utf8]{inputenc}
\usepackage{lmodern}
\usepackage{amsmath,amssymb,amsthm,mathtools}
\usepackage{graphicx,booktabs,xcolor}
\usepackage[hidelinks]{hyperref}
\usepackage[margin=25mm]{geometry}
\theoremstyle{plain}
\newtheorem{theorem}{Theorem}
\newtheorem{lemma}[theorem]{Lemma}
\newtheorem{corollary}[theorem]{Corollary}
\newtheorem{proposition}[theorem]{Proposition}
\theoremstyle{definition}
\newtheorem{definition}[theorem]{Definition}
\newtheorem{example}[theorem]{Example}
\theoremstyle{remark}
\newtheorem*{remark}{Remark}
\setlength{\parindent}{0pt}\setlength{\parskip}{0.6em}
\title{Visual calculus}
\date{}
\begin{document}
\maketitle
Each figure below adds a calculus picture to a \texttt{Figure:} line. Numbers you can move (the
point of tangency, the number of strips, the degree) become sliders next to the graph.

\section*{Tangent lines}

\begin{figure}[htbp]
\centering
\includegraphics[width=0.8\linewidth]{visual-calculus-fig1.png}
\caption{The tangent to a cubic at $x = 1$; drag the point.}
\end{figure}
\section*{Riemann sums}

\begin{figure}[htbp]
\centering
\includegraphics[width=0.8\linewidth]{visual-calculus-fig2.png}
\caption{A left sum with eight strips; it falls short of the curve.}
\end{figure}
\begin{figure}[htbp]
\centering
\includegraphics[width=0.8\linewidth]{visual-calculus-fig3.png}
\caption{A right sum overshoots a rising curve.}
\end{figure}
\begin{figure}[htbp]
\centering
\includegraphics[width=0.8\linewidth]{visual-calculus-fig4.png}
\caption{The midpoint rule is far closer.}
\end{figure}
\begin{figure}[htbp]
\centering
\includegraphics[width=0.8\linewidth]{visual-calculus-fig5.png}
\caption{Trapezoids have slanted tops.}
\end{figure}
\section*{Area between curves}

The curves cross at $0$ and $1$, so the region they enclose is found for you.

\begin{figure}[htbp]
\centering
\includegraphics[width=0.8\linewidth]{visual-calculus-fig6.png}
\caption{The area enclosed by a line and a parabola.}
\end{figure}
\begin{figure}[htbp]
\centering
\includegraphics[width=0.8\linewidth]{visual-calculus-fig7.png}
\caption{Where the curves cross, the shading changes sides.}
\end{figure}
\section*{Solids of revolution}

Drag to turn a solid; the cut shows the disk, washer or shell used.

\begin{figure}[htbp]
\centering
\includegraphics[width=0.8\linewidth]{visual-calculus-fig8.png}
\caption{Disks: turning a square-root curve about the $x$-axis.}
\end{figure}
\begin{figure}[htbp]
\centering
\includegraphics[width=0.8\linewidth]{visual-calculus-fig9.png}
\caption{Washers: the region between a line and a parabola.}
\end{figure}
\begin{figure}[htbp]
\centering
\includegraphics[width=0.8\linewidth]{visual-calculus-fig10.png}
\caption{Shells: turning a parabola about the $y$-axis.}
\end{figure}
\section*{Slope fields and phase portraits}

\begin{figure}[htbp]
\centering
\includegraphics[width=0.8\linewidth]{visual-calculus-fig11.png}
\caption{Solution curves follow the slopes.}
\end{figure}
\begin{figure}[htbp]
\centering
\includegraphics[width=0.8\linewidth]{visual-calculus-fig12.png}
\caption{Logistic growth from two starting values.}
\end{figure}
\begin{figure}[htbp]
\centering
\includegraphics[width=0.8\linewidth]{visual-calculus-fig13.png}
\caption{A centre: trajectories circle the origin.}
\end{figure}
\section*{Gradient, vector fields and contours}

\begin{figure}[htbp]
\centering
\includegraphics[width=0.8\linewidth]{visual-calculus-fig14.png}
\caption{The gradient points straight out, across the level circles.}
\end{figure}
\begin{figure}[htbp]
\centering
\includegraphics[width=0.8\linewidth]{visual-calculus-fig15.png}
\caption{A rotation field.}
\end{figure}
\begin{figure}[htbp]
\centering
\includegraphics[width=0.8\linewidth]{visual-calculus-fig16.png}
\caption{Level curves of a saddle.}
\end{figure}
\section*{Tangent planes}

\begin{figure}[htbp]
\centering
\includegraphics[width=0.8\linewidth]{visual-calculus-fig17.png}
\caption{The plane touching a paraboloid.}
\end{figure}
\section*{Taylor polynomials}

\begin{figure}[htbp]
\centering
\includegraphics[width=0.8\linewidth]{visual-calculus-fig18.png}
\caption{Raise the degree and the polynomial hugs the sine further out.}
\end{figure}
\section*{Newton's method}

\begin{figure}[htbp]
\centering
\includegraphics[width=0.8\linewidth]{visual-calculus-fig19.png}
\caption{Each tangent meets the axis closer to $sqrt 2$.}
\end{figure}
\section*{Data}

A cart rolls down a ramp; its position is measured once a second.

\begin{center}
\begin{tabular}{ll}
\toprule
$t$ (s) & $x$ (m) \\
\midrule
0 & 0.1 \\
1 & 1.2 \\
2 & 3.8 \\
3 & 9.3 \\
4 & 15.7 \\
5 & 25.2 \\
\bottomrule
\end{tabular}
\end{center}
\begin{equation*}
  p(u) \coloneqq  \operatorname{polyreg}(t, x, 2)
\end{equation*}
\begin{figure}[htbp]
\centering
\includegraphics[width=0.8\linewidth]{visual-calculus-fig20.png}
\caption{Measurements with the best parabola through them.}
\end{figure}
\begin{figure}[htbp]
\centering
\includegraphics[width=0.8\linewidth]{visual-calculus-fig21.png}
\caption{How the positions are spread.}
\end{figure}
\begin{figure}[htbp]
\centering
\includegraphics[width=0.8\linewidth]{visual-calculus-fig22.png}
\caption{Median, quartiles and range.}
\end{figure}
\end{document}
